\section*{अध्याय \ref{ch:b1:electronic-effects} --- इलेक्ट्रॉनिक प्रभाव और क्रियाशील मध्यवर्ती}

\begin{solution}{exo:b1:electronic-effects:1}
2-मेथिलब्यूटेन, \ce{CH3-CH(CH3)-CH2-CH3}, में तीनों \ce{CH3} प्राथमिक हैं,
\ce{CH2} द्वितीयक है, और CH तृतीयक। 2-ब्रोमो-2-मेथिलप्रोपेन में तीनों
\ce{CH3} प्राथमिक हैं, केंद्रीय कार्बन तृतीयक। चारों ब्रोमाइडों
\ce{C4H9Br} (1-ब्रोमोब्यूटेन, 2-ब्रोमोब्यूटेन, 1-ब्रोमो-2-मेथिलप्रोपेन,
2-ब्रोमो-2-मेथिलप्रोपेन) में से अंतिम तृतीयक \omterm{def:b1:electronic-effects:intermediates}{कार्बधनायन} देता है, जो
सबसे स्थायी है।
\end{solution}

\begin{solution}{exo:b1:electronic-effects:2}
$10^{4.76 - 2.87} = 10^{1.89} = 78$; $10^{4.76 - 0.52} = 10^{4.24} =
\num{1.7e4}$।
\end{solution}

\begin{solution}{exo:b1:electronic-effects:3}
एथेनोएट: C=O और \ce{C-O^-} दोनों ऑक्सीजनों के बीच अदला-बदली करते हैं।
4-नाइट्रोफ़ीनॉक्साइड: ऑक्सीजन का \omterm{def:b1:lewis-resonance-vsepr:lewis}{एकाकी युग्म} एक C=O बनाता है, वलय के द्वि-आबंध
खिसकते हैं, \ce{NO2} वहन करने वाला कार्बन एक C=N आबंध बनाता है, और एक N=O युग्म
ऑक्सीजन पर जाता है: एक ऐसी संरचना जिसमें वलय के एक सिरे पर C=O, दूसरे पर C=N
और दोनों नाइट्रो ऑक्सीजन ऋणात्मक हैं (एक क्विनॉइड संरचना)। ऐलिल धनायन:
\ce{CH2=CH-CH2+} $\leftrightarrow$ \ce{+CH2-CH=CH2}।
\end{solution}

\begin{solution}{exo:b1:electronic-effects:4}
दूसरा तीर: जल के O--H आबंध से उसके ऑक्सीजन की ओर; उत्पाद \ce{H2O}
(\ce{OH-} से, अब उदासीन) और \ce{OH-} (पूर्व जल का ऑक्सीजन
युग्म रखता है, आवेश $-1$)। अमोनिया के लिए: N के \omterm{def:b1:lewis-resonance-vsepr:lewis}{एकाकी युग्म} से
\ce{H3O+} के एक H की ओर एक तीर, और उस O--H आबंध से ऑक्सीजन की ओर एक तीर: N
$+1$ हो जाता है, O 0।
\end{solution}

\begin{solution}{exo:b1:electronic-effects:5}
एथेनॉल (15.9) $<$ जल (14.0) $<$ फ़ीनॉल (9.99, वलय में विस्थानीकरण) $<$
4-नाइट्रोफ़ीनॉल (7.16, \ce{NO2} आवेश ले लेता है) $<$ एथेनोइक अम्ल (4.76,
दो तुल्य ऑक्सीजन) $<$ ट्राइक्लोरोएथेनोइक अम्ल (0.51, तीन $-I$
क्लोरीन)।
\end{solution}

\begin{solution}{exo:b1:electronic-effects:6}
ऐनिलीन (4.60) $<$ पिरिडीन (5.23) $<$ अमोनिया (9.25) $<$ मेथिलऐमीन
(10.66)। ऐनिलीन में \omterm{def:b1:lewis-resonance-vsepr:lewis}{एकाकी युग्म} वलय के साथ साझा होता है
(C=N$^+$ वाली संरचनाएँ, ऋण आवेश \textit{ऑर्थो} और
\textit{पैरा} पर); N का प्रोटॉनन उस विस्थानीकरण की कीमत पर होगा।
\end{solution}

\begin{solution}{exo:b1:electronic-effects:7}
3-नाइट्रोफ़ीनॉक्साइड में \omterm{def:b1:lewis-resonance-vsepr:resonance}{अनुनादी संरचनाएँ} आवेश को ऑक्सीजन के सापेक्ष
\textit{ऑर्थो} और \textit{पैरा} कार्बनों पर रखती हैं, \ce{NO2} वहन करने वाले कार्बन
पर कभी नहीं: कोई मेसोमेरी सहायता नहीं, इसलिए 4-समावयवी से दुर्बल।
नाइट्रो समूह फिर भी $\sigma$ आबंधों के माध्यम से इलेक्ट्रॉन खींचता है ($-I$),
इसलिए फ़ीनॉल से प्रबल।
\end{solution}

\begin{solution}{exo:b1:electronic-effects:8}
\ce{CH3+} $<$ \ce{CH3CH2+} $<$ \ce{CH2=CH-CH2+} $<$ \ce{C6H5CH2+}
$\approx$ \ce{(CH3)3C+} $<$ \ce{CH3OCH2+}: ऐल्किल समूह इलेक्ट्रॉन देते हैं;
ऐलिल और बेंज़िल आवेश का विस्थानीकरण करते हैं (दो और चार संरचनाएँ);
ऑक्सीजन का \omterm{def:b1:lewis-resonance-vsepr:lewis}{एकाकी युग्म} ऐसी संरचना देता है जिसमें हर जगह अष्टक पूरा है। सूची के
मध्य का क्रम माध्यम पर निर्भर करता है।
\end{solution}

\begin{solution}{exo:b1:electronic-effects:9}
\ce{OH-} (\omterm{def:b1:electronic-effects:nucleophile}{नाभिकरागी}) \ce{CH3Br} (\omterm{def:b1:electronic-effects:nucleophile}{इलेक्ट्रॉनरागी}) के कार्बन पर आक्रमण करता है,
और C--Br युग्म Br पर चला जाता है। \ce{H2O}, \ce{H+} को एक युग्म देता है (यहाँ अम्ल
और क्षारक नाम प्रयुक्त होते हैं)। \ce{CN-} एथेनल के कार्बोनिल कार्बन पर आक्रमण करता है,
और C=O का $\pi$ युग्म ऑक्सीजन पर चला जाता है। \ce{NH3} अपना \omterm{def:b1:lewis-resonance-vsepr:lewis}{एकाकी युग्म}
\ce{BF3} में बोरॉन की रिक्त कक्षक को देता है।
\end{solution}

\begin{solution}{exo:b1:electronic-effects:10}
नहीं: $K_a = 10^{-0.51} = 0.31$, $C$ से बहुत बड़ा है; अम्ल लगभग
पूर्णतः वियोजित है। $x^2/(0.010 - x) = 0.31$ से $x = \qty{9.7e-3}{mol/L}$ मिलता है,
$\mathrm{pH} = 2.01$, \omterm{def:b1:predominant-reaction:dissociation}{वियोजन मात्रा} \qty{97}{\percent}।
\end{solution}

\begin{solution}{exo:b1:electronic-effects:11}
ऐल्किल समूह इलेक्ट्रॉनों को ($+I$) ऐल्कॉक्साइड ऑक्सीजन पर धकेलते हैं और उसके
आवेश को कम स्थायी बनाते हैं, समूह जितने अधिक, उतना अधिक। बड़े ऐल्कॉक्साइड
जल द्वारा कम अच्छी तरह विलायकित भी होते हैं, जो छोटे आवेशित
ऑक्सीजन को ऐल्किल समूहों के बीच दबे ऑक्सीजन से बेहतर स्थायी बनाता है। दोनों
$\mathrm{p}K_a$ बढ़ाते हैं।
\end{solution}

\begin{solution}{exo:b1:electronic-effects:12}
ऐमाइड में नाइट्रोजन का \omterm{def:b1:lewis-resonance-vsepr:lewis}{एकाकी युग्म} C=O के साथ साझा होता है
(\ce{N+=C-O^-} संरचना); उसे प्रोटॉन या \omterm{def:b1:electronic-effects:nucleophile}{इलेक्ट्रॉनरागी} के लिए लेना उस
स्थायित्व को नष्ट कर देगा। \omterm{def:b1:predominant-reaction:strong-acid}{प्रबल अम्ल} में ऐमाइड का प्रोटॉनन
ऑक्सीजन पर होता है, जो आंशिक ऋण आवेश वहन करता है।
\end{solution}

\begin{solution}{pb:b1:electronic-effects:1}
\textbf{1.} $10^{1.89} = 78$; $10^{1.52} = 33$; $10^{0.84} = 6.9$।
\textbf{2.} नहीं: हर अगला क्लोरीन कम सहायता करता है। कार्बोक्सिलेट का
आवेश पहले से ही अच्छी तरह फैला हुआ है, और जैसे-जैसे अम्ल पूर्ण वियोजन के निकट पहुँचता है,
मापन कम यथार्थ होते जाते हैं।
\textbf{3.} हर Cl, C--C और C--O आबंधों के माध्यम से इलेक्ट्रॉन घनत्व आकर्षित करता है
और कार्बोक्सिलेट को स्थायी बनाता है ($-I$)।
\textbf{4.} \omterm{def:b1:electronic-effects:inductive}{प्रेरणिक प्रभाव} दो-तीन आबंधों के भीतर ही क्षीण हो जाता है।
\textbf{5.} दोनों अम्ल समान रूप से प्रबल दिखते हैं (0.52 और 0.51)। जल में,
सामान्य सांद्रताओं पर दोनों लगभग पूर्णतः वियोजित होते हैं
(\cref{ch:b1:predominant-reaction}, समतलन): 0 के पास का $\mathrm{p}K_a$
उस सीमा पर है जहाँ तक जल भेद कर सकता है, और दोनों मानों में
कुछ दशांशों की अनिश्चितता है।
\textbf{6.} एथेनोइक अम्ल: अम्ल रूप (\qty{98}{\percent}); क्लोरोएथेनोइक:
दोनों, ऋणायन थोड़ा प्रभावी ($10^{3 - 2.87} = 1.3$); शेष तीन:
ऋणायन।
\textbf{7.} \ce{CH3-C(=O)-O^-} $\leftrightarrow$ \ce{CH3-C(-O^-)=O}।
\textbf{8.} दो समान C--O आबंध, द्वि- और एकल आबंध के बीच के।
\textbf{9.} आवेश अकेले ऑक्सीजन पर है, और उसे साझा करने के लिए कोई $\pi$ आबंध
नहीं है।
\textbf{10.} $10^{15.9 - 4.76} = 10^{11.1}$, लगभग $10^{11}$।
\textbf{11.} फ़ीनॉक्साइड का आवेश वलय के तीन कार्बनों तक फैलता है, जो ऑक्सीजन से कम
विद्युत-ऋणात्मक हैं: एक लाभ, पर कार्बोक्सिलेट से छोटा:
$\mathrm{p}K_a$ 9.99, 4.76 और 15.9 के बीच।
\textbf{12.} $K = 10^{6.37 - \mathrm{p}K_a}$: एथेनोइक अम्ल $10^{1.61} = 41$,
अनुकूल (हाइड्रोजनकार्बोनेट के आधिक्य के साथ कार्बन डाइऑक्साइड मुक्त होती है);
फ़ीनॉल $10^{-3.62}$ और एथेनॉल $10^{-9.5}$: कोई अभिक्रिया नहीं।
\textbf{13.} 4-नाइट्रोफ़ीनॉल (7.16) $>$ 2-नाइट्रोफ़ीनॉल (7.23) $>$ 3-नाइट्रोफ़ीनॉल
(8.36) $>$ फ़ीनॉल (9.99)।
\textbf{14.} \cref{exo:b1:electronic-effects:3} की क्विनॉइड संरचना,
जिसमें आवेश नाइट्रो समूह के एक ऑक्सीजन पर है।
\textbf{15.} आवेश केवल ऑक्सीजन के सापेक्ष \textit{ऑर्थो} और
\textit{पैरा} कार्बनों तक पहुँचता है; नाइट्रो कार्बन \textit{मेटा} है।
\textbf{16.} \ce{NO2} का $-I$ प्रभाव।
\textbf{17.} 2-नाइट्रोफ़ीनॉल में OH पड़ोसी नाइट्रो समूह के एक ऑक्सीजन के साथ
\omterm{def:b1:intermolecular-forces-solvents:hbond}{हाइड्रोजन आबंध} बनाता है, जो अम्ल रूप को स्थायी बनाता है और
प्रोटॉन को हटाना थोड़ा कठिन कर देता है।
\textbf{18.} ऋणायन अंश $1/(1 + 10^{\mathrm{p}K_a - 7.5})$: 4-नाइट्रो
\qty{69}{\percent}, 2-नाइट्रो \qty{65}{\percent}, 3-नाइट्रो
\qty{12}{\percent}। 4-नाइट्रोफ़ीनॉल का अम्ल और ऋणायन भिन्न ढंग से अवशोषित करते हैं,
ऋणायन दृश्य क्षेत्र में: उसका रंग उसके $\mathrm{p}K_a$ के आसपास प्रकट होता है।
\textbf{19.} $x^2/(0.010 - x) = 10^{-0.52} = 0.30$: $x = \qty{9.7e-3}{mol/L}$,
$\mathrm{pH} = 2.01$: लगभग एक \omterm{def:b1:predominant-reaction:strong-acid}{प्रबल अम्ल}।
\textbf{20.} $\frac12(4.76 + 2.00) = 3.38$।
\textbf{21.} \ce{CF3COOH + CH3COO- -> CF3COO- + CH3COOH}, $K = 10^{4.76 -
0.52} = 10^{4.24}$: \omterm{def:b1:extent-q-and-k:quantitative}{मात्रात्मक}।
\textbf{22.} $K_a(\ce{CF3COOH})/K_a(\ce{CH3COOH}) =$ \textbf{$10^{4.24} =
\num{1.7e4}$}, लगभग दस हज़ार।
\end{solution}
